\begin{document}

\hypertarget{dependency-analysis-and-scheduling-exo-2-report}{%
\section{Dependency Analysis and Scheduling --- Exo 2 (Report)}\label{dependency-analysis-and-scheduling-exo-2-report}}

\textbf{Task:} In21-S8-CS4553 --- Dependency Analysis and Scheduling (Not graded)\\
\textbf{Paper Focus:} \emph{Exo 2: Growing a Scheduling Language} (Appendix A: Scheduling Primitives)

\begin{center}\rule{0.5\linewidth}{0.5pt}\end{center}

\hypertarget{the-separation-of-semantics-and-scheduling-in-exo-2}{%
\subsection{1. Separation of Semantics and Scheduling in Exo 2}\label{the-separation-of-semantics-and-scheduling-in-exo-2}}

Traditional compilers apply a fixed set of optimizations to improve performance while preserving program semantics. When a target architecture provides specialized capabilities (e.g., SIMD/vector extensions, multithreading, or accelerator memory hierarchies), the compiler may attempt to restructure code automatically to exploit them. However, achieving consistently high performance often requires hardware-aware transformations that are difficult to apply safely and predictably in a fully automatic setting.

Exo 2 addresses this by explicitly separating:
\begin{itemize}
  \item \textbf{Semantics / Logic:} a clear reference implementation of the algorithm (the ``spec'' kernel).
  \item \textbf{Scheduling:} a user-defined sequence of verified, semantics-preserving transformations that specialize the implementation for a target architecture (e.g., SIMD, pthreads, CUDA-like memory hierarchies).
\end{itemize}

This design enables aggressive performance tuning while keeping correctness anchored to the reference semantics. Appendix A of the Exo 2 paper lists the scheduling primitives grouped into eight categories.

\begin{center}\rule{0.5\linewidth}{0.5pt}\end{center}

\hypertarget{step-1-scheduling-primitive-categories-appendix-a}{%
\subsection{2. Step 1 --- Scheduling Primitive Categories (Appendix A)}\label{step-1-scheduling-primitive-categories-appendix-a}}

\hypertarget{a.1-loop-transformations}{%
\subsubsection{A.1 Loop Transformations}\label{a.1-loop-transformations}}

\textbf{Primitives:} \texttt{reorder\_loops}, \texttt{divide\_loop}, \texttt{divide\_with\_recompute}, \texttt{mult\_loops},
\texttt{cut\_loop}, \texttt{join\_loops}, \texttt{shift\_loop}, \texttt{fission}, \texttt{remove\_loop},
\texttt{add\_loop}, \texttt{unroll\_loop}

\textbf{Technique:} Reshape the iteration space using loop splitting/tiling, interchange, unrolling, and structural edits (fission/join).

\textbf{Performance impact:} Loop transformations create structures that map efficiently to hardware:
\begin{itemize}
  \item \textbf{SIMD:} fixed-width inner loops (e.g., width 8/16) improve vectorization and intrinsic matching.
  \item \textbf{Threads (pthreads/OpenMP):} outer loops become natural units of parallel work.
  \item \textbf{Accelerators:} blocking/tiling patterns support hierarchical memories (global/shared/register-like).
\end{itemize}

\hypertarget{a.2-code-rearrangement}{%
\subsubsection{A.2 Code Rearrangement}\label{a.2-code-rearrangement}}

\textbf{Primitives:} \texttt{reorder\_stmts}, \texttt{commute\_expr}

\textbf{Technique:} Reorder independent statements and apply safe algebraic commutation.

\textbf{Performance impact:} Reduces artificial dependencies and increases opportunities for vectorization and instruction selection by forming backend-friendly expression patterns.

\hypertarget{a.3-scope-transformations}{%
\subsubsection{A.3 Scope Transformations}\label{a.3-scope-transformations}}

\textbf{Primitives:} \texttt{specialize}, \texttt{fuse}, \texttt{lift\_scope}

\textbf{Technique:} Transform control scopes (loop/if regions) via fusion, specialization, and scope restructuring.

\textbf{Performance impact:}
\begin{itemize}
  \item \textbf{Fusion} improves temporal locality by keeping producer/consumer work in the same loop nest.
  \item \textbf{Specialization} removes unnecessary conditionals when constraints are proven (compile-time simplification).
  \item \textbf{Scope lifting} corrects nesting structures introduced after tiling/splitting.
\end{itemize}

\hypertarget{a.4-multiple-procedures}{%
\subsubsection{A.4 Multiple Procedures}\label{a.4-multiple-procedures}}

\textbf{Primitives:} \texttt{inline}, \texttt{replace}, \texttt{call\_eqv}, \texttt{extract\_subproc}

\textbf{Technique:} Perform inter-procedural scheduling: inline kernels, replace code regions with equivalent optimized procedures, and extract reusable subprocedures.

\textbf{Performance impact:} Enables micro-kernel engineering and instruction selection. In practice, this is how scalar loop bodies can be systematically replaced with optimized vector intrinsics (e.g., AVX2) while preserving semantic equivalence.

\hypertarget{a.5-buffer-transformations}{%
\subsubsection{A.5 Buffer Transformations}\label{a.5-buffer-transformations}}

\textbf{Primitives:} \texttt{lift\_alloc}, \texttt{sink\_alloc}, \texttt{delete\_buffer}, \texttt{reuse\_buffer},
\texttt{resize\_dim}, \texttt{expand\_dim}, \texttt{rearrange\_dim}, \texttt{divide\_dim}, \texttt{mult\_dim},
\texttt{unroll\_buffer}, \texttt{bind\_expr}, \texttt{stage\_mem}

\textbf{Technique:} Control allocation scope and buffer layout: move allocations, reshape dimensions, change indexing/layout, and stage memory windows into temporaries.

\textbf{Performance impact:}
\begin{itemize}
  \item \textbf{Locality/tiling:} \texttt{stage\_mem} implements cache/shared-memory style staging.
  \item \textbf{Footprint reduction:} \texttt{sink\_alloc} with \texttt{resize\_dim} shrinks working sets per tile.
  \item \textbf{Register/SIMD friendliness:} \texttt{expand\_dim}/\texttt{divide\_dim}/\texttt{unroll\_buffer} create small fixed-size temporaries that map to registers or vector lanes.
  \item \textbf{Reuse:} \texttt{reuse\_buffer} reduces allocations when lifetimes do not overlap.
\end{itemize}

\hypertarget{a.6-simplification}{%
\subsubsection{A.6 Simplification}\label{a.6-simplification}}

\textbf{Primitives:} \texttt{simplify}, \texttt{eliminate\_dead\_code}, \texttt{rewrite\_expr}, \texttt{merge\_writes},
\texttt{inline\_window}, \texttt{inline\_assign}

\textbf{Technique:} Post-pass cleanup to remove redundant structures introduced during scheduling.

\textbf{Performance impact:} Produces tighter generated code by removing dead code, simplifying expressions, merging redundant writes, and eliminating unnecessary temporaries.

\hypertarget{a.7-backend-checked-annotations}{%
\subsubsection{A.7 Backend-Checked Annotations}\label{a.7-backend-checked-annotations}}

\textbf{Primitives:} \texttt{set\_memory}, \texttt{set\_precision}, \texttt{parallelize\_loop}, \texttt{set\_window}

\textbf{Technique:} Attach explicit annotations verified by Exo/backends during lowering/code generation.

\textbf{Performance impact:}
\begin{itemize}
  \item \textbf{Parallelism safety:} \texttt{parallelize\_loop} ensures memory safety across loop iterations (no forbidden cross-iteration hazards).
  \item \textbf{SIMD instruction selection:} \texttt{set\_memory} enables memory-aware pattern matching for intrinsic replacement.
  \item \textbf{Precision and interfaces:} \texttt{set\_precision} and \texttt{set\_window} maintain correct typing and window semantics.
\end{itemize}

\hypertarget{a.8-configuration-state}{%
\subsubsection{A.8 Configuration State}\label{a.8-configuration-state}}

\textbf{Primitives:} \texttt{bind\_config}, \texttt{delete\_config}, \texttt{write\_config}

\textbf{Technique:} Schedule configuration writes for stateful accelerators.

\textbf{Performance impact:} Hoists and deduplicates configuration updates so they do not occur redundantly in innermost loops, reducing control/config overhead.

\begin{center}\rule{0.5\linewidth}{0.5pt}\end{center}

\hypertarget{step-2-install-exo-and-run-examples}{%
\subsection{3. Step 2 --- Install Exo and Run Examples}\label{step-2-install-exo-and-run-examples}}

\subsubsection{Install (pip)}
\begin{verbatim}
pip install exo-lang
\end{verbatim}

\subsubsection{Build from source (optional)}
\begin{verbatim}
git clone https://github.com/exo-lang/exo.git
cd exo
git submodule update --init --recursive

python -m venv ~/.venv/exo
source ~/.venv/exo/bin/activate
python -m pip install -U pip setuptools wheel
python -m pip install -r requirements.txt
python -m build .
pip install dist/*.whl
\end{verbatim}

\subsubsection{Run the quizzes}
\begin{verbatim}
cd exo/examples/quiz1 && exocc quiz1.py
cd ../quiz2 && exocc quiz2.py
cd ../quiz3 && exocc quiz3.py
\end{verbatim}

\begin{center}\rule{0.5\linewidth}{0.5pt}\end{center}

\hypertarget{step-3-quiz-scenarios-and-solutions}{%
\subsection{4. Step 3 --- Quiz Scenarios and Solutions}\label{step-3-quiz-scenarios-and-solutions}}

This section summarizes the issues in each quiz, the incorrect scheduling approach, and the corrected solution.

\subsubsection{Quiz 1 --- AVX2 replacement not occurring}

\textbf{Issue:} The schedule forms an 8-wide structure, but the output remains scalar (no AVX2 intrinsics).\\
\textbf{Incorrect approach:} Calling \texttt{replace\_all} before annotating staged buffers with the correct vector memory type.\\
\textbf{Correct solution:} \texttt{replace\_all} is memory-aware. Intermediate buffers such as \texttt{two\_vec}, \texttt{inp\_vec}, and \texttt{out\_vec} must be annotated using \texttt{set\_memory(..., AVX2)} prior to replacement so the intrinsic patterns match.\\
\textbf{Source:} \texttt{solutions/quiz1\_solution.py}

\subsubsection{Quiz 2 --- SchedulingError during loop fission}

\textbf{Issue:} \texttt{SchedulingError}: fission would hide an allocation required by a later use site.\\
\textbf{Incorrect approach:} Performing \texttt{expand\_dim} $\rightarrow$ \texttt{lift\_alloc} $\rightarrow$ \texttt{fission} immediately for each temporary buffer in a single pass.\\
\textbf{Correct solution:} Use a two-pass schedule:
\begin{itemize}
  \item Pass 1: expand and lift \emph{all} temporary allocations to a safe scope.
  \item Pass 2: fission after assignments once allocations are no longer trapped inside the inner loop body.
\end{itemize}
\textbf{Source:} \texttt{solutions/quiz2\_solution.py}

\subsubsection{Quiz 3 --- \texttt{blur\_x} width fails to shrink}

\textbf{Issue:} The allocation height is reduced (to 34), but width remains \texttt{W} instead of shrinking to 256.\\
\textbf{Incorrect approach:} The helper \texttt{get\_loops\_at\_or\_above} failed to include the starting loop cursor (\texttt{xo}), so the allocation was never sunk into the \texttt{xo} tile scope.\\
\textbf{Correct solution:} Initialize the loop list as \texttt{loops=[cursor]} so \texttt{xo} is included. This allows sinking and resizing to correctly reduce the active footprint to \texttt{[34, 256]}.\\
\textbf{Source:} \texttt{solutions/quiz3\_solution.py}

\end{document}